% Part: methods
% Chapter: induction
% Section: structural-induction

\documentclass[../../../include/open-logic-section]{subfiles}

\begin{document}

\olfileid{mth}{ind}{sti}

\olsection{ساختی استقرا}

ہن تک اسیں سارے قدرتی عدداں بارے نتیجے ثابت کرن لئی استقرا ورتیا اے۔ پر ایس ورگا اصول استقرائی طریقے نال تعریف ہوئے سیٹ دے سارے !!{element}s بارے نتیجے سِدھے ای ثابت کرن لئی ورتیا جا سکدا اے۔ اینوں اکثر \emph{ساختی} استقرا آکھدے نیں، کیوں جے ایہ استقرائی طریقے نال تعریف ہوئیاں شےواں دی ساخت اُتے ٹکدا اے۔

عام طور اُتے استقرائی تعریف ایویں دتی جاندی اے۔ (ا) سیٹ دے «مُڈھلے» !!{element}s دی فہرست تے (ب) اوہناں عملاں دی فہرست جیہڑے پرانے رکناں توں سیٹ دے نویں !!{element}s بناؤندے نیں۔ مثلاً ٹھیک ٹرماں دی صورت وچ مُڈھلیاں شےواں اکھر نیں۔ ساڈے کول بس اک عمل اے، ایہ عمل ایہ اے:
\begin{align*}
  o(s_1, s_2) = & [s_1 \circ s_2]
\end{align*}
تسی قدرتی عدداں~$\Nat$ نوں وی استقرائی تعریف نال دتے ہوئے سمجھ سکدے او۔ مُڈھلی شے~$0$ اے تے عمل جانشین فنکشن~$x + 1$ اے۔

استقرائی طریقے نال تعریف ہوئے سیٹ دے سارے رکناں بارے کوئی گل ثابت کرن لئی، یعنی سیٹ دے ہر !!{element} وچ خاصیت~$P$ اے، سانوں ایہ کرنا پئے گا:
\begin{enumerate}
\item ثابت کرو کہ مُڈھلیاں شےواں وچ~$P$ اے۔
\item ثابت کرو کہ ہر عمل~$o$ لئی، جے اوہدے آرگومنٹاں وچ~$P$ اے، تاں نتیجے وچ وی اے۔
\end{enumerate}
مثلاً سارے ٹھیک ٹرماں بارے کوئی گل ثابت کرن لئی، اسیں ثابت کراں گے کہ ایہ سارے اکھراں بارے سچ اے، تے $[s_1 \circ s_2]$ بارے سچ اے بشرطیکہ $s_1$ تے $s_2$ بارے وکھ وکھ سچ ہووے۔

\begin{prop}
  $[$ دی گنتی $]$ دی گنتی برابر اے، کسے وی ٹھیک ٹرم~$t$ وچ۔
\end{prop}

\begin{proof}
اسیں ساختی استقرا ورتدے آں۔ ٹھیک ٹرم استقرائی طریقے نال تعریف ہوئے نیں، اکھر مُڈھلیاں شےواں نیں تے عمل $o$ پرانیاں توں نویں ٹھیک ٹرم بناؤندا اے۔
\begin{enumerate}
\item دعویٰ ہر اکھر لئی سچ اے، کیوں جے اکلّے اکھر وچ $[$ دی گنتی~$0$ اے تے $]$ دی گنتی وی~$0$ اے۔
\item منّ لو $[$ دی گنتی $s_1$ وچ $]$ دی گنتی برابر اے، تے ایہو گل $s_2$ لئی وی سچ اے۔ $[$ دی گنتی $o(s_1, s_2)$ وچ، یعنی $[s_1 \circ s_2]$ وچ، $[$ دیاں $s_1$ تے $s_2$ وچ دیاں گنتیاں دا جوڑ تے نال اک اے۔ $]$ دی گنتی $o(s_1, s_2)$ وچ، $]$ دیاں $s_1$ تے $s_2$ وچ دیاں گنتیاں دا جوڑ تے نال اک اے۔ سو $[$ دی گنتی $o(s_1, s_2)$ وچ $]$ دی گنتی $o(s_1,s_2)$ وچ دے برابر اے۔
\end{enumerate}
\end{proof}

\begin{prob}
  ساختی استقرا نال ثابت کرو کہ کوئی ٹھیک ٹرم~$]$ نال شروع نہیں ہُندا۔
\end{prob}

آؤ ساختی استقرا نال اک ہور ثبوت دیئیے۔ نشاناں دی لڑی~$t$ دا حقیقی مُڈھلا ٹکڑا اوہ لڑی~$s$ اے جیہڑی $t$ نال، کھبّے توں پڑھدیاں، نشان بہ نشان ملدی اے، پر $t$ لمّی اے۔ سو مثلاً $[a \circ {}$، $[a \circ b]$ دا حقیقی مُڈھلا ٹکڑا اے، پر نہ $[b \circ {}$ اے (دوجے نشان اُتے فرق اے) تے نہ $[a
\circ b]$ اے (دوہاں دی لمبائی اکّو اے)۔

\begin{prop}\ollabel{prop:initial}
  کسے ٹھیک ٹرم~$t$ دے ہر حقیقی مُڈھلے ٹکڑے وچ $[$ دی گنتی $]$ دی گنتی توں ودھ اے۔
\end{prop}

\begin{proof}
  ~$t$ اُتے استقرا نال:
  \begin{enumerate}
  \item $t$ اکلّا اکھر اے۔ تاں $t$ دا کوئی حقیقی مُڈھلا ٹکڑا نہیں۔
  \item $t = [s_1 \circ s_2]$ کسے ٹھیک ٹرماں $s_1$ تے~$s_2$ لئی۔ جے $r$، $t$ دا حقیقی مُڈھلا ٹکڑا اے، تاں کئی صورتاں ممکن نیں:
    \begin{enumerate}
    \item $r$ بس $[$ اے۔ تاں $r$ وچ اک $[$،~$]$ نالوں ودھ اے۔
    \item $r$، $[r_1$ اے، جتھے $r_1$،~$s_1$ دا حقیقی مُڈھلا ٹکڑا اے۔ کیوں جے $s_1$ ٹھیک ٹرم اے، استقرائی مفروضے مطابق $r_1$ وچ $[$، $]$ نالوں ودھ نیں، تے ایہو گل $[r_1$ لئی وی سچ اے۔
    \item $r$، $[s_1$ یا $[s_1 \circ {}$ اے۔ پچھلے نتیجے مطابق $[$ تے $]$ دیاں گنتیاں~$s_1$ وچ برابر نیں۔ سو $[$ دی گنتی $[s_1$ یا $[s_1 \circ {}$ وچ،~$]$ دی گنتی نالوں اک ودھ اے۔
    \item $r$، $[s_1 \circ r_2$ اے، جتھے $r_2$،~$s_2$ دا حقیقی مُڈھلا ٹکڑا اے۔ استقرائی مفروضے مطابق $r_2$ وچ $[$، $]$ نالوں ودھ نیں۔ پچھلے نتیجے مطابق $[$ تے $]$ دیاں گنتیاں~$s_1$ وچ برابر نیں۔ سو $[$ دی گنتی $[s_1 \circ r_2$ وچ،~$]$ دی گنتی توں ودھ اے۔
    \item $r$، $[s_1 \circ s_2$ اے۔ پچھلے نتیجے مطابق $[$ تے $]$ دیاں گنتیاں $s_1$ وچ برابر نیں، تے~$s_2$ لئی وی ایہو گل اے۔ سو اک $[$، $[s_1 \circ s_2$ وچ،~$]$ دی گنتی نالوں ودھ اے۔
    \end{enumerate}
  \end{enumerate}
\end{proof}

\end{document}
